\documentclass[11pt,a4paper]{article}

% --- Paquetes requeridos ---
\usepackage[utf8]{inputenc}
\usepackage[spanish,es-nodecimaldot]{babel}
\usepackage[margin=2.5cm]{geometry}
\usepackage{amsmath,amsfonts,amssymb}
\usepackage{booktabs}
\usepackage{tabularx}
\usepackage{microtype}
\usepackage{hyperref}
\usepackage{float}
\usepackage{comment}

% --- Configuración de hipervínculos ---
\hypersetup{
	colorlinks=true,
	linkcolor=blue,
	urlcolor=blue,
	citecolor=blue
}

\title{\textbf{Taller de Física: Actividad Exploratoria y Cinemática}}
\author{\textbf{}}
%\date{\today}

\begin{document}
	
	\maketitle
	
	\section{Actividad Exploratoria: Conceptos Fundamentales}
	
	\subsection{Definiciones de conceptos}
	
	\begin{itemize}
		\item \textbf{Caída libre:} Movimiento rectilíneo, hacia abajo, que ocurre cuando se deja caer un objeto desde cierta altura.
		\item \textbf{Movimiento parabólico:} Tipo de movimiento que se genera al lanzar algo con cierta inclinación, por ejemplo al disparar un cañón.
		\item \textbf{Trayectoria:} Es el recorrido de un cuerpo o la ``huella'' que deja al moverse.
		\item \textbf{Gravedad:} Es la aceleración que aparece cerca de objetos con masa; el valor en la superficie terrestre es aproximadamente $9{,}8\,\text{m/s}^2$.
		\item \textbf{Velocidad:} Es el cambio de posición con respecto al tiempo.
		\item \textbf{Altura:} Es la distancia que hay desde un punto de referencia hasta otro situado verticalmente por encima de él.
	\end{itemize}
	
	\subsection{Para los dibujos}
	Se proponen las siguientes situaciones para los dibujos:
	
	\begin{enumerate}
		\item \textbf{Caída libre:} Silueta de un edificio alto y una mano soltando una esfera hacia el suelo.
		\item \textbf{Movimiento parabólico:} Un jugador de baloncesto o un cañón disparando un proyectil con un ángulo de elevación, trazando una trayectoria en forma de parábola hacia el objetivo.
		\item \textbf{Trayectoria:} Un avión en vuelo dejando una estela punteada continua que indica el camino espacial recorrido.
		\item \textbf{Gravedad:} El planeta Tierra con un objeto descendiendo hacia su superficie, acompañado por un vector de aceleración vertical $\vec{g}$ orientado hacia el centro de masa terrestre.
		\item \textbf{Velocidad:} Un automóvil en línea recta sobre una vía con la leyenda de $t=0$ y el mismo carro más adelante con la leyenda de $t=5$ o cualesquiera números diferentes.
		\item \textbf{Altura:} Una torre o estructura vertical con una cota de medida ($h$ o $y$) que conecta la base con la cúspide.
	\end{enumerate}
	
	\subsection{Comparación con los Conceptos de google}
	
	\begin{table}[H]
		\centering
		\small
		\begin{tabularx}{\textwidth}{l p{3.8cm} p{4.2cm} X}
			\toprule
			\textbf{Concepto} & \textbf{Definición del Cuaderno} & \textbf{Definición Formal} & \textbf{Comparación} \\
			\midrule
			\textbf{Caída libre} & Movimiento vertical descendente al soltar un objeto. & Movimiento de un cuerpo bajo la influencia exclusiva de un campo gravitatorio, sin resistencia del medio. & Coinciden en el descenso, pero la formal habla estrictamente sobre la fricción con el aire. \\
			\addlinespace
			\textbf{Mov. parabólico} & Movimiento originado al lanzar algo con inclinación. & Movimiento bidimensional resultante de componer un MRU horizontal con un MRUA vertical bajo la gravedad. & La mía habla del lanzamiento a cierto ángulo; la formal explica la composición de movimientos. \\
			\addlinespace
			\textbf{Trayectoria} & Es el recorrido o ``huella'' que deja al moverse un cuerpo. & Lugar geométrico de las sucesivas posiciones que ocupa un punto material en movimiento. & Se parecen entre sí pero la formal está en palabras de definiciones matemáticas, en lugar de hablar de una "huella" \\
			\addlinespace
			\textbf{Gravedad} & Es la aceleración que aparece cerca de objetos con masa, ($9{,}8\,\text{m/s}^2$ en la Tierra). & Fuerza de atracción gravitatoria mutua o aceleración producida por un campo gravitacional sobre una masa testigo. & Ambas hablan de la gravedad como consecuencia de la masa pero la versión formal enfatiza la interacción mutua entre masas. \\
			\addlinespace
			\textbf{Velocidad} & Cambio de posición respecto al tiempo. & Magnitud vectorial que define el desplazamiento instantáneo por unidad de tiempo: $\vec{v} = \frac{d\vec{r}}{dt}$. & La definición informal menciona la tasa temporal y la formal especifica su carácter vectorial y direccional. \\
			\addlinespace
			\textbf{Altura} & Distancia que hay desde un punto hasta otro más arriba. & Distancia o coordenada vertical medida sobre un eje ortogonal a partir de un nivel de referencia fijado. & Coinciden, la formal explicita la necesidad de fijar un marco de referencia. \\
			\bottomrule
		\end{tabularx}
		\caption{Contraste entre definiciones intuitivas y formalismo físico.}
		\label{tab:comparacion_conceptos}
	\end{table}
	
	\vspace{1em}
	\section{Ensayo Crítico: El Mito de la Caída Libre y la Física Real}
	
	\subsection*{Introducción}
	El análisis del ingeniero aeroespacial pone en discusión algo que es ampliamente difundido: la afirmación de que todos los objetos caen igual de rápido sin importar su masa. Si bien esto suele transmitirse como una regla en la educación básica, su validez depende estrictamente de condiciones ideales que no existen en el entorno en donde vivimos.
	
	\subsection*{La Hipótesis del Vacío vs la Resistencia Aerodinámica}
	El principio de Galileo Galilei, comprobado por la misión Apolo XV en la Luna al dejar caer una pluma y un martillo, establece que en el \emph{vacío} la aceleración gravitacional es idéntica para cualquier cuerpo:
	\begin{equation}
		g \approx 9{,}81\,\text{m/s}^2 \qquad\text{en la tierra}
	\end{equation}
	En este escenario, al no existir atmósfera, la masa del objeto no afecta cuanto tiempo demora en caer.
	
	Sin embargo, en la tierra existe atmósfera, lo que genera una \textbf{fuerza de resistencia} ($F_r$) que se opone al movimiento descendente:
	\begin{equation}
		F_r = \frac{1}{2}\,\rho\,v^2\,A\,C_d
	\end{equation}
	donde $\rho$ representa la densidad del aire, $v$ la velocidad instantánea, $C_d$ el coeficiente aerodinámico de arrastre y $A$ el área de la sección transversal expuesta.
	
	A medida que un cuerpo cae y acelera, la fuerza de rozamiento se incrementa proporcionalmente a la velocidad al cuadrado ($v^2$) hasta igualar la fuerza del peso ($W = m \cdot g$). Cuando ambas fuerzas se equilibran ($\sum F_y = 0$), la aceleración neta se anula y el cuerpo alcanza su denominada \textbf{velocidad terminal}. %($v_t$):
	
	Esto significa que la masa ($m$) sí interviene de forma decisiva en la velocidad terminal ya que dependiendo de la masa el objeto tiene un peso mayor o no. Dos personas que caen juntas en un salto tándem duplican la masa pero mantienen un área frontal similar, alcanzando una velocidad terminal sensiblemente mayor a la de un paracaidista individual. De igual manera, el salto de Félix Baumgartner desde más de $39.000\,\text{m}$ de altitud permitió superar la barrera del sonido debido a que la densidad atmosférica ($\rho$) a esa altura es menor a la de la superficie terrestre.
	
	\subsection*{Conclusión}
	El estudio de la caída libre demuestra que la física real no debe reducirse a fórmulas memorizadas bajo condiciones ideales. Comprender el papel de la resistencia aerodinámica y de la velocidad límite permite superar los mitos y reconocer la complejidad con la que interactúan las fuerzas en nuestro entorno cotidiano.
	
\end{document}